%%%%%%%%%%  scan idx 186  |  folio 173  %%%%%%%%%%
inversible sur $S$ est trivial. Il est immédiat d'autre part que $p$ induit un isomorphisme
\[
p^*:\Gamma(\mathbb E^2_k,\mathcal O_{\mathbb E^2_k})\longrightarrow\Gamma(S,\mathcal O_S).
\]
Par suite, toute section globale de $\mathcal O_S$ qui s'annule en $O_1$ s'annule aussi en $O_2$, donc les $S_f$, pour $f\in\Gamma(S,\mathcal O_S)$ ne forment pas une base de la topologie de $S$.

\textbf{Lemme 2.2.8.} Soient $S$ un schéma divisoriel (2.2.5), $(L_i)_{i\in I}$ une famille ample de $\mathcal O_S$-Modules inversibles. Notons $L(S)$ la sous-catégorie additive strictement pleine de $\Qcoh(S)$ (2.2.1) engendrée par les $L_i^{\otimes-n}$ pour $(i,n)\in I\times\mathbb N^+$. Soit $E\in\Ob\mathrm D(\Qcoh(S))$.
\begin{enumerate}[label=\alph*)]
\item Les catégories $L(S)$ et $\Loclib(S)$ (2.0) sont quasi-relevables dans $\Qcoh(S)$ (I 1.2) et quasi-stables par noyau d'épimorphisme (I 1.2).
\item Soit $n\in\mathbb Z$. Pour que $E$ soit $n$-pseudo-cohérent (resp. pseudo-cohérent) (2.2.1) il faut et il suffit que $E$ soit $n$-pseudo-cohérent (resp. pseudo-cohérent) relativement à $L(S)$ ou $\Loclib(S)$ (I 2.3). Les foncteurs canoniques
\[
\mathrm K^{-}(L(S))/\mathrm K^{-,\varnothing}(L(S))\longrightarrow \mathrm D(\Qcoh(S))
\]
et
\[
\mathrm K^{-}(\Loclib(S))/\mathrm K^{-,\varnothing}(\Loclib(S))\longrightarrow \mathrm D(\Qcoh(S))
\]
sont pleinement fidèles et ont pour image essentielle $\mathrm D^{-}(\Qcoh(S))_{\coh}$ (2.2.1).
\item Soit $[m,n]$ un intervalle de $\mathbb Z$. Pour que $E$ soit d'amplitude parfaite contenue dans $[m,n]$ (2.2.1), il faut et il suffit que $E$ soit isomorphe, dans $\mathrm D(\Qcoh(S))$, à un complexe $F$ tel que $F^i=0$ pour $i\notin[m,n]$, $F^i\in\Ob L(S)$ pour $i>m$, $F^m$ étant localement libre de type fini. Le foncteur canonique
\[
\mathrm K^b(\Loclib(S))/\mathrm K^{b,\varnothing}(\Loclib(S))\longrightarrow \mathrm D(\Qcoh(S))
\]
